\section{What the Operator Listens To}\label{sec:listening}

The input of a classical operator is fixed: a parent and a few
parameters. The input of the LLM operator is a prompt of thousands of
tokens: the task, the parent code, its score, its error trace and the
mutation instruction. AAD methods engineer this prompt on the
assumption that the model uses all of it.
To see which parts of the prompt the model uses, we rewrote the
feedback of 192 parent prompts and generated four children per variant
with common random numbers (4,180 evaluated children,
Table~\ref{tab:counterfactual}), and traced attribution
graphs~\cite{ameisen2025circuits} from the first edited token back to
the prompt. For a crashed parent, deleting the error trace raises the
child's failure rate from 43\% to 66\% (+0.22, CI 0.14 to 0.33), and a
generic error message recovers only part of it; in the attribution
graphs, the error text carries 61\% of the influence on the edit token
in 14 of 16 pairs. For a working parent, rewriting its score from 0.01
to 0.90, or deleting the feedback, changes neither the child's failure
rate nor its fitness, and the internal fitness readout moves by only
$\pm0.04$. Steering along the fitness representation changes how much
of the parent is rewritten, not the outcome.

The replication on Qwen3-14B and Gemma-3-12B
(Table~\ref{tab:cf-models}) confirms the main point and
qualifies the rest. No model uses the written score (Gemma's failures
rise slightly with a low score, +0.08, one of many effects tested).
Qwen3-14B uses the error trace as Qwen3-8B does (+0.24 failures
without it), but Gemma does not: it reads the trace (it repeats the
same exception more often without it), yet its failure rate does not
change. The effect of steering along fitness does not replicate. For
these models, the score in the feedback is decorative; the parent code
and, model-dependently, the error trace drive the edit.

For the design of AAD methods this has two consequences. Feedback is
useful where it carries information the model can act on: the error
trace tells it what to fix, while the score tells it nothing it uses,
so the selection pressure of the loop lies in which parents are shown,
not in what is said about them. And the prompt is a coarse control
channel: parts of it are ignored, and an explicit instruction, as we
show in Sec.~\ref{sec:loop}, moves the code but also imprints its
wording on it. This motivates controlling the operator from inside.
